\chapter{Conclusion} \label{chap:conclusion}

This thesis explores the use of automatic summarization as more than a tool of information compression. Given advances in language generation technology, I explore the possibility of using summarization methodology as a tool for information communication. In particular, I focus on building methods for the generation and evaluation of scientific summaries, with a focus on the needs and preferences of the end user. This thesis is rooted in the core belief that scientific discovery should be accessible, and that accessibility starts but does not end at open-access. From a high-school student who is considering a college major, to a concerned parent wondering what medical decisions to make for their child, to a researcher trying to make sense of our world, everyone benefits from an understanding of the latest in science. Yet each of these individuals is just that - an individual with individual needs. In order to truly make science accessible to all, we need to build systems that take into account the needs and use-cases of each person, creating summaries that are tailored to each person's background and informational needs.

\Cref{part1} of this thesis focuses on designing systems that follow user-indicated attributes of controllability. The core idea of this section is that a user may have different requirements for different use-cases, and that allowing users to indicate their desired output allows for more useful summaries.~\Cref{chap:chap-4} considers the attribute of structure. A number of documents, such as blogs, posters, and slides, are often structured summaries of an underlying scientific document. Our framework is a two-step process that first builds a structured representation of the knowledge present in a scientific document, then leverages that representation to generate a summary, following a user-specified structure. Then in~\Cref{chap:chap-5}, we consider the axis of readability. We aim to create a ``dial'' of readability, that allows the user to indicate on a spectrum the desired level of technicality in a summary, rather than a simple binary choice between lay and technical. In this chapter, we use control vectors - a representation engineering method that extracts a vector that can be applied at inference time to control the level of readability of a generation. This method is training-free and low-data, making it a significantly lower compute method compared to past work. We find that this method, in combination with our proposed requirements-based framework for data selection, gives us better controllability over strong baselines. Although the majority of this thesis considers scientific accessibility from the perspective of the readability and utility of the output summaries,~\Cref{part1} focuses on building methods that improve the performance of smaller models, thus increasing the accessibility via decreasing the compute cost. In envisioning systems that make science accessible to everyone, this includes working towards systems in which the cost to use them is not overly prohibitive. 

\Cref{part2} of this thesis focuses on designing methods for the evaluation of user-centric summarization systems. The majority of past work on summarization evaluation focuses on evaluating the overall quality of a summary. However, in building methods for user-centric generation, we must also develop methods that allow us to better measure whether or not a summary meets the needs of that user.~\Cref{chap:chap-5} first looks at how we evaluate readability. We find that the most popular metrics used to measure readability suffer from definitional inconsistency and measurement error, resulting in low agreement with human judgments of readability. In other words, the most popular metrics of readability were designed in an educational context, and do not generalize to the task of readable summarization. We find that LLMs are better judges of readability, able to measure deeper features of readability, such as explanation of background concepts. We recommend that future work in plain language summarization use a multi-faceted approach to evaluation, leveraging language models, better performing traditional metrics, and qualitative analysis. In~\Cref{chap:chap-6}, I consider the case of low annotation budget for evaluation. While past work has shown that LLMs are better judges when provided with a gold reference, realistically gold references are not always available. This is especially true in our setting; collecting gold references for highly personalized summaries is more expensive and difficult than collecting generally high quality summaries, as the axes of personalization greatly increase the number of required summaries. Thus we explore the setting in which related but non-exact references are available. For example, if we have access to related summaries in a different style, or even cheaply-acquired, machine-generated summaries from other models. We refer to this as the ``reference-adjusted'' setting, a middle ground between reference-based and reference-free. We show that this setting consistently outperforms the reference-free setting, and approaches the reference-based setting. Finally, in~\Cref{chap:chap-7}, we introduce the axis of evaluation called \textit{informational satisfaction}. In other words, does a summary meet the user's informational needs? I show that the majority of evaluation metrics do not pass simple perturbations tests regarding the informational content of a summary. I then conduct a human evaluation in which I collect preference data about informational satisfaction and find that current methods, including LLM-as-judge metrics, do not accurately capture human preference on information satisfaction. Overall, this part makes steps towards user-centric summarization evaluation through both new methodology and analysis pointing to gaps in current metrics.

\section{Resources Created}

Through the process of conducting research for this thesis, I have introduced a number of artifacts for the community, which I hope will support future research in this area, or even simply use of my systems. Below I compile a list of the resources I've created, followed by a brief description of each:
\begin{itemize}[label={\faGithub},noitemsep]
    \item \repo{https://github.com/microsoft/knowledge-centric-templatic-views}{microsoft/knowledge-centric-templatic-views} {$\rightarrow$} Contains the implementation for the method and evaluation metric introduced in~\Cref{chap:chap-3}.
    \item \repo{https://github.com/JHU-CLSP/control-summ}{JHU-CLSP/control-summ} {$\rightarrow$} Contains the implementation for the control vector method and baselines described in~\Cref{chap:chap-4}.
    \item\repo{https://github.com/JHU-CLSP/eval-the-eval-readability}{JHU-CLSP/eval-the-eval-readability}{$\rightarrow$} Contains the analysis and evaluation code described in~\Cref{chap:chap-5}.
    \item \repo{https://github.com/JHU-CLSP/ref-adj-eval}{JHU-CLSP/ref-adj-eval} {$\rightarrow$} Contains the analysis code for the evaluation framework introduced in~\Cref{chap:chap-6}.
    \item \repo{https://github.com/isabelcachola/persona-eval-metric}{isabelcachola/persona-eval-metric} {$\rightarrow$} Contains the implementation for the evaluation metric and baselines described in~\Cref{chap:chap-7}.
    \item \repo{https://github.com/isabelcachola/persona-eval-annotation}{isabelcachola/persona-eval-annotation} {$\rightarrow$} Contains the implementation for the annotation platform used in~\Cref{chap:chap-7}.
\end{itemize}

\section{Ethics}
This thesis involves the use of Large Language Models for generation and evaluation. LLM generated documents can potentially generate copy-righted material~\cite{carlini21extracting}, personally-identifiable information~\cite{Lukas2023AnalyzingLO}, or factually incorrect text~\cite{Manakul2023SelfCheckGPTZB}. The use of LLMs to generate documents may violate some academic dishonesty policies~\cite{Zdravkova2023IntegrationOL}.
Although the majority of research in controllable generation focuses on general controllability~\cite{Keskar2019CTRLAC,Chen2024BenchmarkingLL,yang-klein-2021-fudge} and AI Safety features~\cite{zou2025representationengineeringtopdownapproach,rimsky-etal-2024-steering}, the same methods could be used for nefarious tasks, such as generating purposely dishonest language~\cite{BARMAN2024100545}. In using LLMs for evaluation, LLMs are subject to bias~\cite{Venkit2024AnAO,Stureborg2024LargeLM}. Additionally, the use of LLMs contributes to the environmental footprint of our field~\cite{Schwartz2019GreenA}.  

My thesis focuses on designing methods with the ultimate goal of making scientific discovery more accessible. However, we have to contend with the possible harms, no matter how well-intentioned the source research may be. In this thesis, we have made efforts to mitigate possible harms when appropriate (e.g. focusing on smaller models that require less compute). The potential risks of conducting this research are common among many areas of research within the Artificial Intelligence and Natural Language Processing fields. However, they have become more immediately relevant in recent years as consumer use of the systems we build has only increased. More people and more companies are incorporating LLMs into their everyday workflows. As with prior technological advances of the past, this sharp increase in usage has impacted society in both positive and negative ways. The use of LLMs can cause harms both unintentionally (e.g. misinformation via generating hallucinated text) and intentionally (e.g. purposely generating harmful text). These are problems that we must contend with as a research community and as citizens of a technologically dependent world. Although I believe the potential benefits of my research outweigh the possible harms, this is a conversation I will continue to engage in with my fellow researchers for the rest of my career. It is our responsibility to not only build technology that improves people's lives, but also mitigate the potential harms of such technology. 

\section{Limitations}\label{sec:thesis-limitations}

I divide the limitations of this thesis into 2 categories: limitations as a result of experimental design and limitations of the larger impact of this body of work.

\paragraph{Limitations in Experimental Design.} First is the concern of generalizability; every chapter draws its experiments from English-language scientific text. Whether the methods and findings extend to other languages or to domains such as law, clinical notes, or educational material remains an open question. The generative methods presented are confined to textual content, leaving multi-modal generation and evaluation to future work. A second thread of limitations concerns the depth of mechanistic understanding behind the proposed methods. The control vector experiments, for instance, establish downstream effects on readability without verifying that the vectors themselves encode readability as a latent representation, and a fuller account of what such vectors capture would require further investigation. The experiments presented in this thesis generally rely on relatively small, open-source judges. Since stronger models tend to yield better results overall, the reported improvements likely represent a lower bound on what larger or proprietary models could achieve, although generalization of our specific methods is not explored. Finally, the possibility of test data contamination in pretraining corpora cannot be ruled out, though this concern applies symmetrically to LLM-baselines and therefore does not undermine comparative claims.

\paragraph{Limitations in Broader Impact.} The stated goal of this thesis is to build systems that make science more accessible by focusing on the needs and wants of the end user. While I believe that this thesis makes a step in this direction, there is more work to be done to make this vision a reality. First, the generative methods presented in~\Cref{part1} focus on 2 axes of controllability, structure and readability. These are obviously not the only axes of import. This provides 2 primary limitations of the generative methods introduced. First, they focus on specific personalization axes, and do not address the question of generalized personalization. To achieve a fully personalized experience, we need to answer a number of questions. What axes of personalization are important to users? What is the best interaction design? How do users consider the trade-off between privacy and personalization? How much personalization should be user-indicated versus machine-inferred? Each of these questions is not answered by this thesis, and they are avenues of future work. The second limitation of the generative methods is that they are disjoint. We introduce a method for improving the structure of summaries and a method for controlling the readability levels. This is limited in that they are specifically designed, and are not guaranteed to generalize, or even be applicable, to other axes of personalization. Similar limitations can be said of the evaluation methodologies introduced in~\Cref{part2}. We focus on the axes of information needs and readability, but do not provide a unified evaluation method that can be used for any personalization task. Such unification of personalized summarization methodologies is a distinct limitation of my work, and a good avenue for future work.


\section{Future Work}
% \epigraph{Not everything important is measurable, and \\
    % not everything measurable is important.}{Elliot Eisner}
{
    \small
    \textit{Not everything important is measurable, 
    and not everything measurable is important.} \\
    \textsc{-- Elliot Eisner}
}\\

The future of summarization isn't about better compression. It's about better understanding of the people we're summarizing for.

For most of its history, the field has treated summarization as an information bottleneck problem~\cite{knight2002summarization, clarke2008global}. Take a long text, produce a shorter one, preserve what matters. Recent directions like plain language and query-focused summarization complicate this picture, but most modern systems still bake in rigid assumptions about who the reader is and what they need~\cite{August2024KnowYA,Zhang2023BenchmarkingLL}. My own work on plain language summarization made this concrete. The same summary can be perfectly faithful to a source document and still fail a reader who lacks the background knowledge to parse it. ``What matters'' isn't a property of the text, it's a property of the reader.

This points to a deeper methodological problem. NLP operates in cycles. Build a benchmark, solve the benchmark, repeat. This cycle worked when failures were easy to quantify (ungrammatical outputs, off-task generations, hallucinated facts). Benchmarks still catch these, and they should. But the failures that matter most for real users are harder to see in aggregate metrics. For example, a summary that's technically correct but assumes expertise the reader doesn't have, an answer that's fluent but subtly misaligned with what was asked, or a system that performs well on average but fails the people who needed it most. We've optimized for what's measurable rather than what's valuable.

So what comes next? I'd argue the interesting questions are no longer about capability but about fit. Where does AI genuinely improve people's lives, and where is human judgment irreplaceable? What does a system look like when it's designed around a specific user in a specific context, rather than a benchmark distribution? ``Stop chasing leaderboards'' is a tired refrain, so let me offer something concrete: the methods we need already exist, just not in our field. Human-computer interaction has spent decades developing tools for understanding users, contextual inquiry, participatory design, longitudinal field studies, and NLP has only scratched the surface. A summarization system co-designed with the clinicians, patients, or students who will actually use it will look different from one tuned on ROUGE, and we won't know how different until we try.

The pace of the field will date any specific proposals I might offer quickly, so I'll end with the orientation instead of a list. Build for people, with people, and let the systems follow.